% Luc Destombe --- licence CC BY-NC-SA 4.0
% https://creativecommons.org/licenses/by-nc-sa/4.0/deed.fr
% Fichier autonome : compiler avec pdflatex, deux fois (signets, nombre de pages).
\documentclass[11pt,a4paper]{article}
\makeatletter
\newif\iflycee@unecolonne
\lycee@unecolonnetrue
\newif\iflycee@palatino
\lycee@palatinotrue

\RequirePackage[utf8]{inputenc}
\RequirePackage[T1]{fontenc}
\RequirePackage[french]{babel}
\RequirePackage{etoolbox}

\newcommand{\ShorthandsOff}{\@ifundefined{shorthandoff}{}{\shorthandoff{;:!?}}}
\newcommand{\ShorthandsOn}{\@ifundefined{shorthandon}{}{\shorthandon{;:!?}}}
\ShorthandsOff
\AtBeginDocument{\ShorthandsOn}
\AtBeginEnvironment{tikzpicture}{\ShorthandsOff}

\RequirePackage{geometry}
\RequirePackage{fancyhdr}
\RequirePackage{lastpage}
\RequirePackage{multicol}
\RequirePackage{array}
\RequirePackage{enumitem}
\RequirePackage{xcolor}
\RequirePackage{needspace}

\RequirePackage{amsmath,amssymb}
\RequirePackage{mathpazo}
\renewcommand{\ttdefault}{lmtt}
\DeclareMathSizes{10.95}{10.95}{8}{5.5}
\begingroup\catcode`\_=\active \catcode`\^=\active
\gdef_#1{\sb{\scriptscriptstyle #1}}
\gdef^#1{\sp{\scriptscriptstyle\lycee@moinsepais #1}}
\endgroup
\newcommand\lycee@moins{\mathbin{%
  \setbox\z@\hbox{$\m@th\scriptscriptstyle\lycee@moinsorig$}%
  \dimen@=.55\wd\z@ \advance\dimen@-.5pt
  \hbox to.75\wd\z@{\hss$\m@th\scriptscriptstyle\vcenter{\hbox{%
    \tikz\draw[line width=.5pt,line cap=round](0,0)--(\the\dimen@,0);}}$\hss}}}
\begingroup\lccode`\~=`\- \lowercase{\endgroup\let~}\lycee@moins
\newcommand\lycee@moinsepais{\mathcode`\-="8000 }
\AtBeginDocument{\mathchardef\lycee@moinsorig=\mathcode`\-\relax}
\AtBeginDocument{\catcode`\_=12 \mathcode`\_="8000
  \catcode`\^=12 \mathcode`\^="8000 }
\newcommand\lycee@exposants{%
  \fontdimen13\textfont2=.48\fontdimen6\textfont2
  \fontdimen14\textfont2=.48\fontdimen6\textfont2
  \fontdimen15\textfont2=.48\fontdimen6\textfont2
  \fontdimen13\scriptfont2=.48\fontdimen6\scriptfont2
  \fontdimen14\scriptfont2=.48\fontdimen6\scriptfont2
  \fontdimen15\scriptfont2=.48\fontdimen6\scriptfont2}
\every@math@size\expandafter{\the\every@math@size\lycee@exposants}
\newlength{\retraitcalculs}
\setlength{\retraitcalculs}{2em}

\RequirePackage{tikz}
\usetikzlibrary{arrows.meta,calc,babel,shapes.misc}
\RequirePackage[most]{tcolorbox}
\tcbuselibrary{skins,breakable}

\definecolor{coulDef}{HTML}{1F4E79}
\definecolor{coulProp}{HTML}{7030A0}
\definecolor{coulRem}{HTML}{C55A11}
\definecolor{coulEx}{HTML}{2E7D32}
\definecolor{coulExo}{HTML}{B03A2E}
\definecolor{coulPart}{HTML}{1F4E79}
\definecolor{coulMeth}{HTML}{1B6E4F}
\definecolor{coulCourbe}{HTML}{0B5ED7}
\definecolor{coulCourbeB}{HTML}{C00000}
\definecolor{coulCourbeC}{HTML}{2E7D32}
\definecolor{coulGrille}{HTML}{8C8C8C}
\definecolor{coulRep}{HTML}{1B6E4F}
\definecolor{coulBar}{HTML}{2E7D32}

\newcommand{\R}{\mathbb{R}}
\newcommand{\Rs}{\mathbb{R}^{*}}
\newcommand{\N}{\mathbb{N}}
\newcommand{\Ns}{\mathbb{N}^{*}}
\newcommand{\Z}{\mathbb{Z}}
\newcommand{\Q}{\mathbb{Q}}
\renewcommand{\le}{\leqslant}
\renewcommand{\ge}{\geqslant}
\renewcommand{\approx}{\simeq}
\newcommand{\vect}[1]{\overrightarrow{#1}}
\newcommand{\norme}[1]{\left\lVert #1 \right\rVert}
\newcommand{\intervalle}[2]{\left[#1\,;#2\right]}
\RequirePackage{cancel}
\renewcommand{\CancelColor}{\color{coulExo}}
\newcommand{\fois}[1]{\times{\color{coulExo}#1}}
\newcommand{\barre}[1]{\cancel{#1}}
\newlength{\lycee@eqh}\newlength{\lycee@eqd}
\newsavebox{\lycee@eqa}\newsavebox{\lycee@eqb}\newsavebox{\lycee@eqc}
\newcommand{\lycee@eqcol}[1]{\hbox{$\renewcommand{\arraystretch}{1.35}%
  \begin{array}[t]{@{}>{\displaystyle}r@{}>{\displaystyle{}}l@{}}#1\end{array}$}}
\newcommand{\lycee@eqmesure}[1]{%
  \ifdim\ht#1>\lycee@eqh\setlength{\lycee@eqh}{\ht#1}\fi
  \ifdim\dp#1>\lycee@eqd\setlength{\lycee@eqd}{\dp#1}\fi}
\newcommand{\lycee@eqtrait}{\hspace{1.2em}\tikz[baseline=0pt]\draw[dash pattern=on 2pt off 2pt,
  line width=0.5pt] (0,\lycee@eqh) -- (0,-\lycee@eqd);\hspace{1.2em}}
\newcommand{\eqproduit}[2]{\par\smallskip\noindent C'est une équation produit :
  $A\times B=0 \iff A=0$ ou $B=0$.\par\nopagebreak\smallskip
  \begingroup\setlength{\lycee@eqh}{0pt}\setlength{\lycee@eqd}{0pt}%
  \setbox\lycee@eqa=\lycee@eqcol{#1}\setbox\lycee@eqb=\lycee@eqcol{#2}%
  \lycee@eqmesure\lycee@eqa\lycee@eqmesure\lycee@eqb%
  \noindent\hspace*{\retraitcalculs}\usebox\lycee@eqa\lycee@eqtrait\usebox\lycee@eqb\endgroup\par\smallskip}
\newcommand{\eqproduittrois}[3]{\par\smallskip\noindent C'est une équation produit :
  $A\times B\times C=0 \iff A=0$ ou $B=0$ ou $C=0$.\par\nopagebreak\smallskip
  \begingroup\setlength{\lycee@eqh}{0pt}\setlength{\lycee@eqd}{0pt}%
  \setbox\lycee@eqa=\lycee@eqcol{#1}\setbox\lycee@eqb=\lycee@eqcol{#2}\setbox\lycee@eqc=\lycee@eqcol{#3}%
  \lycee@eqmesure\lycee@eqa\lycee@eqmesure\lycee@eqb\lycee@eqmesure\lycee@eqc%
  \noindent\hspace*{\retraitcalculs}\usebox\lycee@eqa\lycee@eqtrait\usebox\lycee@eqb\lycee@eqtrait\usebox\lycee@eqc\endgroup
  \par\smallskip}

\tikzset{grille/.style={coulGrille,line width=0.4pt}}
\newcommand{\quadrillage}[4]{\draw[grille,step=1] (#1,#2) grid (#3,#4);}
\tikzset{
  croix/.style={cross out,draw,line width=0.6pt,inner sep=0pt,minimum size=4pt},
  croixdroite/.style={croix,rotate=45,line width=0.8pt},
}
\newcommand{\pt}[3]{\path (#1) node[croix]{} node[#2,font=\small,inner sep=2.5pt]{$#3$};}

\newcommand{\lycee@repere}[4]{%
  \draw[grille,step=1] (#1,#2) grid (#3,#4);
  \draw[-{Stealth[length=2mm]},thick] (#1,0)--({#3+0.4},0) node[below left=-1pt]{$x$};
  \draw[-{Stealth[length=2mm]},thick] (0,#2)--(0,{#4+0.4}) node[below left=-1pt]{$y$};
  \node[below left,font=\scriptsize,inner sep=1.5pt] at (0,0) {$0$};}
\newcommand{\lycee@gradx}[1]{\foreach \k/\l in {#1}{%
  \draw (\k,0.07)--(\k,-0.07) node[below,font=\scriptsize,inner sep=1.5pt]{$\l$};}}
\newcommand{\lycee@grady}[1]{\foreach \k/\l in {#1}{%
  \draw (0.07,\k)--(-0.07,\k) node[left,font=\scriptsize,inner sep=1.5pt]{$\l$};}}
\newcommand{\lycee@intff}[2]{\left[#1\,;#2\right]}
\newcommand{\lycee@intfo}[2]{\left[#1\,;#2\right[}
\newcommand{\lycee@intof}[2]{\left]#1\,;#2\right]}
\newcommand{\lycee@intoo}[2]{\left]#1\,;#2\right[}
\newcommand{\lycee@ens}[1]{\left\{#1\right\}}
\AtBeginDocument{%
  \providecommand{\repere}{\lycee@repere}%
  \providecommand{\gradx}{\lycee@gradx}%
  \providecommand{\grady}{\lycee@grady}%
  \providecommand{\Cf}{\mathcal{C}\sb{\scriptscriptstyle f}}%
  \providecommand{\Cg}{\mathcal{C}\sb{\scriptscriptstyle g}}%
  \providecommand{\intff}{\lycee@intff}%
  \providecommand{\intfo}{\lycee@intfo}%
  \providecommand{\intof}{\lycee@intof}%
  \providecommand{\intoo}{\lycee@intoo}%
  \providecommand{\ens}{\lycee@ens}}

\newlist{questions}{enumerate}{3}
\setlist[questions,1]{label=\arabic*\textdegree),leftmargin=*,itemsep=2pt,topsep=3pt}
\setlist[questions,2]{label=\alph*),leftmargin=*,itemsep=1pt,topsep=2pt}
\setlist[questions,3]{label=\roman*),leftmargin=*,itemsep=1pt,topsep=1pt}
\setlist[itemize]{label=\textbullet,leftmargin=*,itemsep=2pt,topsep=3pt}

\newcommand{\besoinplace}[1]{\par\penalty\@M\begingroup
  \dimen@=#1\relax\dimen@ii=\pagegoal\advance\dimen@ii-\pagetotal
  \ifdim\dimen@>\dimen@ii\ifdim\dimen@ii>\z@\vfil\fi\break\fi\endgroup}

\newcommand{\lycee@pied}{\thepage/\pageref{LastPage}}
\newcommand{\entete}[2]{%
  \pagestyle{fancy}\fancyhf{}%
  \fancyhead[L]{\footnotesize\color{coulPart}\textbf{#1}}%
  \fancyhead[R]{\footnotesize\color{coulPart}\textbf{#2}}%
  \fancyfoot[C]{\footnotesize\lycee@pied}%
  \renewcommand{\headrulewidth}{0.4pt}%
  \renewcommand{\headrule}{\hbox to\headwidth{%
    \color{coulPart!40}\leaders\hrule height \headrulewidth\hfill}}}

\RequirePackage{ccicons}
\newcommand{\lycee@licence}{{\scriptsize\color{gray}%
  \copyright\ Luc Destombe --- \ccbyncsa\ CC BY-NC-SA 4.0}}
\newif\iflycee@licence \lycee@licencetrue
\newcommand{\sanslicence}{\lycee@licencefalse}
\AtBeginDocument{\iflycee@licence\AddToHookNext{shipout/foreground}{%
  \put(0,0){\rlap{\kern\dimexpr1in+\hoffset+\oddsidemargin+\textwidth\relax
    \llap{\raisebox{-\dimexpr1in+\voffset+\topmargin+\headheight+\headsep
      +\textheight+\footskip\relax}{\lycee@licence}}}}}\fi}

\newcommand{\lycee@site}{faire-des-maths.fr}
\newcommand{\lycee@version}{{\scriptsize\color{gray}\lycee@site
  \ifdefined\lyceeversion\ --- Version~\lyceeversion\fi}}
\AtBeginDocument{\iflycee@licence\AddToHookNext{shipout/foreground}{%
  \put(0,0){\rlap{\kern\dimexpr1in+\hoffset+\oddsidemargin\relax
    \raisebox{-\dimexpr1in+\voffset+\topmargin+\headheight+\headsep
      +\textheight+\footskip\relax}{\lycee@version}}}}\fi}

\newcommand{\formulesserrees}{%
  \setlength{\abovedisplayskip}{3pt}\setlength{\belowdisplayskip}{3pt}%
  \setlength{\abovedisplayshortskip}{2pt}\setlength{\belowdisplayshortskip}{3pt}}

\setlength{\parindent}{0pt}

\RequirePackage[implicit=false,hidelinks,pdfencoding=auto,psdextra,bookmarksopen,
  bookmarksopenlevel=1,pdfcreator={},pdfproducer={}]{hyperref}
\RequirePackage{bookmark}
\hypersetup{pdfauthor={Luc Destombe},pdfkeywords={CC BY-NC-SA 4.0}}
\newcounter{lycee@signet}
\newcommand{\signet}[3][]{\stepcounter{lycee@signet}%
  \edef\lycee@dest{\if\relax\detokenize{#1}\relax
    signet.\arabic{lycee@signet}\else#1\fi}%
  \hypertarget{\lycee@dest}{}\bookmark[level=#2,dest=\lycee@dest]{#3}}
\pdfstringdefDisableCommands{%
  \def\R{R}\def\N{N}\def\Z{Z}\def\Q{Q}\def\vect#1{#1}\def\Cf{Cf}\def\Cg{Cg}\def\fois#1{×#1}%
  \def\barre#1{#1}\def\intervalle#1#2{[#1;#2]}\def\intff#1#2{[#1;#2]}%
  \def\textsuperscript#1{#1}\def\dfrac#1#2{#1/#2}\def\frac#1#2{#1/#2}%
  \def\vec#1{#1⃗}}%
\RequirePackage[official]{eurosym}

\tcbset{exercice corrige/.style={
  enhanced,breakable,before upper=\raggedright,
  colback=white,colframe=coulExo,
  boxrule=0.8pt,arc=2pt,
  left=6pt,right=6pt,top=8pt,bottom=5pt,
  before skip=9pt,after skip=4pt,
  coltitle=white,fonttitle=\bfseries\small,
  attach boxed title to top left={xshift=8pt,yshift=-\tcboxedtitleheight/2},
  boxed title style={colback=coulExo,arc=2pt,boxrule=0pt}}}
\newcounter{exo}
\newtcolorbox{exercice}[1][]{exercice corrige,
  phantom={\signet[exercice-\arabic{exo}]{2}{Exercice \arabic{exo}}},
  title={Exercice \theexo\ifx\relax#1\relax\else{} --- #1\fi}}
\newenvironment{exo}[1][\relax]{\refstepcounter{exo}\begin{exercice}[#1]}{\end{exercice}}

\geometry{top=1.8cm,bottom=1cm,left=1cm,right=1cm,headheight=14pt,footskip=0.6cm}
\renewcommand{\lycee@pied}{\thepage}

\renewcommand{\theexo}{\ifnum\value{exo}<10 0\fi\arabic{exo}}
\tcbset{exercice corrige/.style={
  enhanced,breakable,before upper=\raggedright,
  colback=white,colframe=coulExo,
  boxrule=0.9pt,arc=2pt,
  left=8pt,right=8pt,top=8pt,bottom=6pt,
  before skip=8pt,after skip=6pt,
  coltitle=white,fonttitle=\bfseries,
  attach boxed title to top left={xshift=8pt,yshift=-\tcboxedtitleheight/2},
  boxed title style={colback=coulExo,arc=2pt,boxrule=0pt}}}

\newcommand{\entetefiche}[2]{%
  \begin{center}
    \begin{tcolorbox}[enhanced,width=0.92\linewidth,colback=white,colframe=coulPart,
        boxrule=1.4pt,arc=3pt,halign=center,top=6pt,bottom=6pt]
      {\large\bfseries\color{coulPart} CORRIGÉ --- FICHE D'EXERCICES}\\[3pt]
      {\normalsize\bfseries #1}\\[2pt]
      {\normalsize #2}
    \end{tcolorbox}
  \end{center}}

\newcounter{partie}
\newcommand{\partie}[1]{%
  \stepcounter{partie}%
  \besoinplace{8\baselineskip}\vspace{4pt}%
  \begin{tcolorbox}[enhanced,colback=coulPart!8,colframe=coulPart,boxrule=0pt,
      leftrule=3pt,arc=0pt,left=6pt,right=4pt,top=3pt,bottom=3pt,
      before skip=6pt,after skip=2pt]
    \bfseries\color{coulPart}\Roman{partie} --- #1\signet{1}{\Roman{partie} --- #1}
  \end{tcolorbox}}
\newcommand{\partiesansnum}[1]{%
  \besoinplace{8\baselineskip}\vspace{4pt}%
  \begin{tcolorbox}[enhanced,colback=coulPart!8,colframe=coulPart,boxrule=0pt,
      leftrule=3pt,arc=0pt,left=6pt,right=4pt,top=3pt,bottom=3pt,
      before skip=6pt,after skip=2pt]
    \bfseries\color{coulPart}#1\signet{1}{#1}
  \end{tcolorbox}}

\setlist[questions,1]{label=\arabic*\textdegree),leftmargin=*,itemsep=3pt,topsep=3pt}
\setlist[questions,2]{label=\alph*),leftmargin=*,itemsep=2pt,topsep=2pt}
\setlist[questions,3]{label=\roman*),leftmargin=*,itemsep=1pt,topsep=1pt}
\newlist{expr}{enumerate}{1}
\setlist[expr]{label={\Alph*\ $=$},leftmargin=*,itemsep=2pt,topsep=3pt}

\newcommand{\res}[1]{{\color{coulRep}#1}}
\newcommand{\rem}[1]{\par\smallskip{\small\itshape\color{black!60}$\rightarrow$ #1}}
\newenvironment{calculs}{\par\smallskip\noindent$\displaystyle\begin{aligned}[t]}%
  {\end{aligned}$\par\smallskip}
\newcommand{\justif}[1]{\quad\text{\small\itshape\color{black!60}#1}}
\newcommand{\cas}[1]{\par\smallskip\noindent\textbf{\color{coulExo}#1}\vadjust{\nobreak}}
\newcommand{\explique}[2]{\par\smallskip\noindent
  \begin{minipage}[t]{0.6\linewidth}\raggedright #1\end{minipage}\hfill
  \begin{minipage}[t]{0.37\linewidth}\raggedright\leavevmode
    {\small\itshape\color{black!60}#2}\end{minipage}\par}

\setlength{\parskip}{2pt}

\makeatother
\entete{Chapitre 13 --- Systèmes d'équations --- corrigé}{Seconde}
\setlength{\parindent}{0pt}
\setlength{\jot}{5pt}

\let\casligne\cas
\renewcommand{\cas}[1]{\casligne{#1}\par\nopagebreak}

\ShorthandsOff
\tikzset{
  droite/.style={coulCourbe,line width=1pt},
  droiteB/.style={coulCourbeB,line width=1pt},
  aide/.style={dashed,black!50,line width=0.5pt},
}
\ShorthandsOn
\newcommand{\sys}[1]{\left\{\begin{aligned}#1\end{aligned}\right.}
\newcommand{\coord}[2]{\begin{pmatrix}#1\\#2\end{pmatrix}}
\newcommand{\sysligne}[1]{\begin{center}${\sys{#1}}$\end{center}}

\newcommand{\cotefig}[3][0.32]{\par\vspace{-0.55\baselineskip}\noindent
  \begin{minipage}[t]{\dimexpr0.97\linewidth-#1\linewidth\relax}\raggedright #2\end{minipage}\hfill
  \begin{minipage}[t]{#1\linewidth}\vspace{0pt}\centering #3\end{minipage}\par}

\begin{document}

\entetefiche{Chapitre 13 --- Systèmes d'équations}{Seconde}

{\small Le repère est orthonormé. Dans chaque système, $(L_{1})$ est la première
équation et $(L_{2})$ la deuxième. Pour la droite d'équation $ax+by+c=0$, le vecteur
$\vec{u}\coord{-b}{a}$ est un vecteur directeur ; pour $y=mx+p$, c'est
$\vec{u}\coord{1}{m}$.\par}

\partie{Systèmes par substitution}

\begin{exo}[Être solution d'un système]
\explique{Un couple est solution quand il vérifie \textbf{les deux}
équations.}{définition 1}
\cas{$(3\,;2)$}
\explique{$3\times 3+2=11$ et $3-2\times 2=-1$ : le couple $(3\,;2)$ \res{est
solution} du système.}{les deux équations sont vérifiées}
\cas{$(4\,;-1)$}
\explique{$3\times 4+(-1)=11$, mais $4-2\times(-1)=6$, et non $-1$ : le couple
$(4\,;-1)$ \res{n'est pas solution} du système.}{il ne vérifie que $(L_{1})$}
\cas{$(1\,;1)$}
\explique{$3\times 1+1=4$, et non $11$ : le couple $(1\,;1)$ \res{n'est pas
solution} du système.}{inutile de tester $(L_{2})$}
\rem{Erreur fréquente : s'arrêter à la première équation. Un couple qui vérifie
$(L_{1})$ seulement n'est pas solution.}
\end{exo}

\begin{exo}[Une inconnue déjà isolée]
\explique{Une équation donne déjà une inconnue : on la remplace dans l'autre
équation.}{méthode 1}
\cas{a)}
\begin{calculs}
2x+y &= 12 &&\justif{$(L_{2})$}\\
2x+(x+3) &= 12 &&\justif{on remplace $y$ par $x+3$}\\
3x+3 &= 12 &&\\
3x &= 9 &&\justif{on enlève $3$}\\
x &= 3 &&\justif{on divise par $3$}
\end{calculs}
\begin{calculs}
y &= 3+3 &&\justif{dans $(L_{1})$ : $y=x+3$}\\
y &= 6 &&
\end{calculs}
\explique{$S=\res{\{(3\,;6)\}}$.}{contrôle dans $(L_{2})$ : $2\times 3+6=12$}
\cas{b)}
\begin{calculs}
3x-4y &= 1 &&\justif{$(L_{2})$}\\
3(2y-1)-4y &= 1 &&\justif{on remplace $x$ par $2y-1$}\\
6y-3-4y &= 1 &&\justif{on développe}\\
2y-3 &= 1 &&\justif{on réduit}\\
2y &= 4 &&\justif{on ajoute $3$}\\
y &= 2 &&\justif{on divise par $2$}
\end{calculs}
\begin{calculs}
x &= 2\times 2-1 &&\justif{dans $(L_{1})$ : $x=2y-1$}\\
x &= 3 &&
\end{calculs}
\explique{$S=\res{\{(3\,;2)\}}$.}{contrôle dans $(L_{2})$ : $9-8=1$}
\cas{c)}
\begin{calculs}
4x-y &= 7 &&\justif{$(L_{2})$}\\
4x-(-2x+5) &= 7 &&\justif{on remplace $y$ par $-2x+5$, entre parenthèses}\\
4x+2x-5 &= 7 &&\justif{le « $-$ » devant la parenthèse change les signes}\\
6x &= 12 &&\justif{on ajoute $5$}\\
x &= 2 &&\justif{on divise par $6$}
\end{calculs}
\begin{calculs}
y &= -2\times 2+5 &&\justif{dans $(L_{1})$}\\
y &= 1 &&
\end{calculs}
\explique{$S=\res{\{(2\,;1)\}}$.}{contrôle dans $(L_{2})$ : $8-1=7$}
\rem{Erreur fréquente : oublier les parenthèses en remplaçant, et écrire $4x-2x+5$.}
\end{exo}

\begin{exo}[Isoler, puis remplacer]
\explique{On isole l'inconnue de coefficient $1$ ou $-1$, puis on la remplace dans
l'autre équation.}{méthode 1}
\cas{a)}
\begin{calculs}
x &= 5-2y &&\justif{on isole $x$ dans $(L_{1})$ : on enlève $2y$}
\end{calculs}
\begin{calculs}
3x+5y &= 13 &&\justif{$(L_{2})$}\\
3(5-2y)+5y &= 13 &&\justif{on remplace $x$ par $5-2y$}\\
15-6y+5y &= 13 &&\justif{on développe}\\
15-y &= 13 &&\justif{on réduit}\\
-y &= -2 &&\justif{on enlève $15$}\\
y &= 2 &&\justif{on change les signes}
\end{calculs}
\begin{calculs}
x &= 5-2\times 2 &&\justif{$x=5-2y$}\\
x &= 1 &&
\end{calculs}
\explique{$S=\res{\{(1\,;2)\}}$.}{contrôle dans $(L_{2})$ : $3+10=13$}
\cas{b)}
\begin{calculs}
y &= 7-2x &&\justif{on isole $y$ dans $(L_{1})$}
\end{calculs}
\begin{calculs}
5x+3y &= 18 &&\justif{$(L_{2})$}\\
5x+3(7-2x) &= 18 &&\justif{on remplace $y$ par $7-2x$}\\
5x+21-6x &= 18 &&\justif{on développe}\\
-x+21 &= 18 &&\justif{on réduit}\\
-x &= -3 &&\justif{on enlève $21$}\\
x &= 3 &&\justif{on change les signes}
\end{calculs}
\begin{calculs}
y &= 7-2\times 3 &&\justif{$y=7-2x$}\\
y &= 1 &&
\end{calculs}
\explique{$S=\res{\{(3\,;1)\}}$.}{contrôle dans $(L_{2})$ : $15+3=18$}
\cas{c)}
\begin{calculs}
3x-y &= 4 &&\justif{$(L_{1})$}\\
-y &= 4-3x &&\justif{on enlève $3x$}\\
y &= 3x-4 &&\justif{on change tous les signes}
\end{calculs}
\begin{calculs}
5x-2y &= 6 &&\justif{$(L_{2})$}\\
5x-2(3x-4) &= 6 &&\justif{on remplace $y$ par $3x-4$}\\
5x-6x+8 &= 6 &&\justif{on développe : $-2\times(-4)=8$}\\
-x+8 &= 6 &&\justif{on réduit}\\
-x &= -2 &&\justif{on enlève $8$}\\
x &= 2 &&\justif{on change les signes}
\end{calculs}
\begin{calculs}
y &= 3\times 2-4 &&\justif{$y=3x-4$}\\
y &= 2 &&
\end{calculs}
\explique{$S=\res{\{(2\,;2)\}}$.}{contrôle dans $(L_{2})$ : $10-4=6$}
\cas{d)}
\begin{calculs}
-x+4y &= 6 &&\justif{$(L_{1})$}\\
-x &= 6-4y &&\justif{on enlève $4y$}\\
x &= 4y-6 &&\justif{on change tous les signes}
\end{calculs}
\begin{calculs}
2x-3y &= 3 &&\justif{$(L_{2})$}\\
2(4y-6)-3y &= 3 &&\justif{on remplace $x$ par $4y-6$}\\
8y-12-3y &= 3 &&\justif{on développe}\\
5y-12 &= 3 &&\justif{on réduit}\\
5y &= 15 &&\justif{on ajoute $12$}\\
y &= 3 &&\justif{on divise par $5$}
\end{calculs}
\begin{calculs}
x &= 4\times 3-6 &&\justif{$x=4y-6$}\\
x &= 6 &&
\end{calculs}
\explique{$S=\res{\{(6\,;3)\}}$.}{contrôle dans $(L_{2})$ : $12-9=3$}
\end{exo}

\begin{exo}[Des solutions en fractions]
\cas{a)}
\begin{calculs}
x &= 2-3y &&\justif{on isole $x$ dans $(L_{1})$}
\end{calculs}
\begin{calculs}
2x-3y &= 1 &&\justif{$(L_{2})$}\\
2(2-3y)-3y &= 1 &&\justif{on remplace $x$ par $2-3y$}\\
4-6y-3y &= 1 &&\justif{on développe}\\
-9y &= -3 &&\justif{on enlève $4$}\\
y &= \frac{-3}{-9} &&\justif{on divise par $-9$}\\
y &= \frac{1}{3} &&\justif{on simplifie par $-3$}
\end{calculs}
\begin{calculs}
x &= 2-3\times\frac{1}{3} &&\justif{$x=2-3y$}\\
x &= 2-1 &&\justif{$3\times\frac13=1$}\\
x &= 1 &&
\end{calculs}
\explique{$S=\res{\left\{\left(1\,;\frac{1}{3}\right)\right\}}$.}{contrôle dans
$(L_{2})$ : $2-3\times\frac13=2-1=1$}
\cas{b)}
\begin{calculs}
5x-2y &= 2 &&\justif{$(L_{2})$}\\
5x-2(x+1) &= 2 &&\justif{on remplace $y$ par $x+1$}\\
5x-2x-2 &= 2 &&\justif{on développe}\\
3x &= 4 &&\justif{on réduit et on ajoute $2$}\\
x &= \frac{4}{3} &&\justif{on divise par $3$}
\end{calculs}
\begin{calculs}
y &= \frac{4}{3}+1 &&\justif{dans $(L_{1})$ : $y=x+1$}\\
y &= \frac{4}{3}+\frac{1\fois{3}}{1\fois{3}} &&\justif{même dénominateur}\\
y &= \frac{7}{3} &&
\end{calculs}
\explique{$S=\res{\left\{\left(\frac{4}{3}\,;\frac{7}{3}\right)\right\}}$.}{contrôle
dans $(L_{2})$ : $\frac{20}{3}-\frac{14}{3}=\frac63=2$}
\end{exo}

\partie{Systèmes par combinaison}

\begin{exo}[Des coefficients déjà opposés]
\explique{Une inconnue a des coefficients opposés : on additionne $(L_{1})$ et
$(L_{2})$, elle disparaît.}{méthode 2}
\cas{a)}
\begin{calculs}
(2x+3y)+(4x-3y) &= 12+6 &&\justif{$(L_{1})+(L_{2})$}\\
6x &= 18 &&\justif{les $y$ s'éliminent : $3y-3y=0$}\\
x &= 3 &&\justif{on divise par $6$}
\end{calculs}
\begin{calculs}
2\times 3+3y &= 12 &&\justif{on remplace $x$ par $3$ dans $(L_{1})$}\\
6+3y &= 12 &&\\
3y &= 6 &&\justif{on enlève $6$}\\
y &= 2 &&\justif{on divise par $3$}
\end{calculs}
\explique{$S=\res{\{(3\,;2)\}}$.}{contrôle dans $(L_{2})$ : $12-6=6$}
\cas{b)}
\begin{calculs}
(-5x+2y)+(5x+4y) &= -1+13 &&\justif{$(L_{1})+(L_{2})$}\\
6y &= 12 &&\justif{les $x$ s'éliminent}\\
y &= 2 &&\justif{on divise par $6$}
\end{calculs}
\begin{calculs}
-5x+2\times 2 &= -1 &&\justif{dans $(L_{1})$}\\
-5x+4 &= -1 &&\\
-5x &= -5 &&\justif{on enlève $4$}\\
x &= 1 &&\justif{on divise par $-5$}
\end{calculs}
\explique{$S=\res{\{(1\,;2)\}}$.}{contrôle dans $(L_{2})$ : $5+8=13$}
\cas{c)}
\begin{calculs}
(3x+4y)+(-3x+2y) &= 1+5 &&\justif{$(L_{1})+(L_{2})$}\\
6y &= 6 &&\justif{les $x$ s'éliminent}\\
y &= 1 &&\justif{on divise par $6$}
\end{calculs}
\begin{calculs}
3x+4\times 1 &= 1 &&\justif{dans $(L_{1})$}\\
3x &= -3 &&\justif{on enlève $4$}\\
x &= -1 &&\justif{on divise par $3$}
\end{calculs}
\explique{$S=\res{\{(-1\,;1)\}}$.}{contrôle dans $(L_{2})$ : $3+2=5$}
\end{exo}

\begin{exo}[Multiplier une équation]
\explique{On multiplie une équation pour que les coefficients d'une inconnue
deviennent opposés, puis on additionne.}{méthode 2}
\cas{a)}
\explique{En $x$, $3$ et $6$ : on multiplie $(L_{1})$ par $-2$.}{$3\times(-2)=-6$,
opposé de $6$}
\sysligne{-6x-4y &= -16 && (L_{1})\times(-2)\\ 6x+5y &= 17 && (L_{2})}
\begin{calculs}
(-6x-4y)+(6x+5y) &= -16+17 &&\justif{on additionne}\\
y &= 1 &&\justif{les $x$ s'éliminent}
\end{calculs}
\begin{calculs}
3x+2\times 1 &= 8 &&\justif{dans $(L_{1})$}\\
3x &= 6 &&\justif{on enlève $2$}\\
x &= 2 &&\justif{on divise par $3$}
\end{calculs}
\explique{$S=\res{\{(2\,;1)\}}$.}{contrôle dans $(L_{2})$ : $12+5=17$}
\cas{b)}
\explique{En $x$, $4$ et $2$ : on multiplie $(L_{2})$ par $-2$.}{$2\times(-2)=-4$,
opposé de $4$}
\sysligne{4x+3y &= 5 && (L_{1})\\ -4x-10y &= -26 && (L_{2})\times(-2)}
\begin{calculs}
(4x+3y)+(-4x-10y) &= 5-26 &&\justif{on additionne}\\
-7y &= -21 &&\justif{les $x$ s'éliminent}\\
y &= 3 &&\justif{on divise par $-7$}
\end{calculs}
\begin{calculs}
4x+3\times 3 &= 5 &&\justif{dans $(L_{1})$}\\
4x &= -4 &&\justif{on enlève $9$}\\
x &= -1 &&\justif{on divise par $4$}
\end{calculs}
\explique{$S=\res{\{(-1\,;3)\}}$.}{contrôle dans $(L_{2})$ : $-2+15=13$}
\cas{c)}
\explique{En $y$, $2$ et $-4$ : on multiplie $(L_{1})$ par $2$.}{$2\times 2=4$, opposé
de $-4$}
\sysligne{10x+4y &= 2 && (L_{1})\times 2\\ 3x-4y &= 11 && (L_{2})}
\begin{calculs}
(10x+4y)+(3x-4y) &= 2+11 &&\justif{on additionne}\\
13x &= 13 &&\justif{les $y$ s'éliminent}\\
x &= 1 &&\justif{on divise par $13$}
\end{calculs}
\begin{calculs}
5\times 1+2y &= 1 &&\justif{dans $(L_{1})$}\\
2y &= -4 &&\justif{on enlève $5$}\\
y &= -2 &&\justif{on divise par $2$}
\end{calculs}
\explique{$S=\res{\{(1\,;-2)\}}$.}{contrôle dans $(L_{2})$ : $3+8=11$}
\rem{Ici, éliminer $y$ est plus simple : une seule multiplication suffit. On regarde
les deux inconnues avant de choisir.}
\end{exo}

\begin{exo}[Multiplier les deux équations]
\explique{On multiplie chaque équation par le coefficient de l'autre, avec un signe
« $-$ » sur l'une des deux : les coefficients deviennent opposés.}{méthode 2 c)}
\cas{a)}
\explique{En $x$, les coefficients sont $3$ et $5$ ; on multiplie $(L_{1})$ par $5$ et
$(L_{2})$ par $-3$.}{$15$ et
$-15$}
\sysligne{15x+25y &= -5 && (L_{1})\times 5\\ -15x-9y &= -27 && (L_{2})\times(-3)}
\begin{calculs}
(15x+25y)+(-15x-9y) &= -5-27 &&\justif{on additionne}\\
16y &= -32 &&\justif{les $x$ s'éliminent}\\
y &= -2 &&\justif{on divise par $16$}
\end{calculs}
\begin{calculs}
3x+5\times(-2) &= -1 &&\justif{dans $(L_{1})$}\\
3x-10 &= -1 &&\\
3x &= 9 &&\justif{on ajoute $10$}\\
x &= 3 &&\justif{on divise par $3$}
\end{calculs}
\explique{$S=\res{\{(3\,;-2)\}}$.}{contrôle dans $(L_{2})$ : $15-6=9$}
\cas{b)}
\explique{En $x$, les coefficients sont $2$ et $5$ ; on multiplie $(L_{1})$ par $5$ et
$(L_{2})$ par $-2$.}{$10$ et
$-10$}
\sysligne{10x+15y &= 20 && (L_{1})\times 5\\ -10x-8y &= -6 && (L_{2})\times(-2)}
\begin{calculs}
(10x+15y)+(-10x-8y) &= 20-6 &&\justif{on additionne}\\
7y &= 14 &&\justif{les $x$ s'éliminent}\\
y &= 2 &&\justif{on divise par $7$}
\end{calculs}
\begin{calculs}
2x+3\times 2 &= 4 &&\justif{dans $(L_{1})$}\\
2x &= -2 &&\justif{on enlève $6$}\\
x &= -1 &&\justif{on divise par $2$}
\end{calculs}
\explique{$S=\res{\{(-1\,;2)\}}$.}{contrôle dans $(L_{2})$ : $-5+8=3$}
\cas{c)}
\explique{En $y$, les coefficients sont $-3$ et $2$ ; on multiplie $(L_{1})$ par $2$ et
$(L_{2})$ par $3$.}{$-6$ et $6$,
déjà de signes contraires}
\sysligne{8x-6y &= 26 && (L_{1})\times 2\\ 9x+6y &= 42 && (L_{2})\times 3}
\begin{calculs}
(8x-6y)+(9x+6y) &= 26+42 &&\justif{on additionne}\\
17x &= 68 &&\justif{les $y$ s'éliminent}\\
x &= 4 &&\justif{on divise par $17$}
\end{calculs}
\begin{calculs}
4\times 4-3y &= 13 &&\justif{dans $(L_{1})$}\\
-3y &= -3 &&\justif{on enlève $16$}\\
y &= 1 &&\justif{on divise par $-3$}
\end{calculs}
\explique{$S=\res{\{(4\,;1)\}}$.}{contrôle dans $(L_{2})$ : $12+2=14$}
\end{exo}

\begin{exo}[Choisir sa méthode]
\cas{a)}
\explique{Dans $(L_{1})$, $y$ a pour coefficient $1$ : la substitution est la plus
rapide.}{remarque « Choisir sa méthode »}
\begin{calculs}
y &= x-7 &&\justif{on isole $y$ dans $(L_{1})$ : on ajoute $x$}
\end{calculs}
\begin{calculs}
2x+y &= 8 &&\justif{$(L_{2})$}\\
2x+(x-7) &= 8 &&\justif{on remplace $y$ par $x-7$}\\
3x &= 15 &&\justif{on réduit et on ajoute $7$}\\
x &= 5 &&\justif{on divise par $3$}
\end{calculs}
\explique{$y=5-7=-2$, donc $S=\res{\{(5\,;-2)\}}$.}{contrôle dans $(L_{2})$ :
$10-2=8$}
\rem{La combinaison marche aussi : les coefficients de $y$ sont égaux, on multiplie
$(L_{1})$ par $-1$, puis on additionne.}
\cas{b)}
\explique{$(L_{1})$ donne déjà $x$ : la substitution est la plus rapide.}{}
\begin{calculs}
2x-5y &= 4 &&\justif{$(L_{2})$}\\
2(3y+1)-5y &= 4 &&\justif{on remplace $x$ par $3y+1$}\\
6y+2-5y &= 4 &&\justif{on développe}\\
y &= 2 &&\justif{on réduit et on enlève $2$}
\end{calculs}
\explique{$x=3\times 2+1=7$, donc $S=\res{\{(7\,;2)\}}$.}{contrôle dans $(L_{2})$ :
$14-10=4$}
\cas{c)}
\explique{Aucun coefficient ne vaut $1$ ou $-1$ : la combinaison évite les
fractions.}{$(L_{1})\times 3$ et $(L_{2})\times(-2)$}
\sysligne{12x+9y &= 30 && (L_{1})\times 3\\ -12x+10y &= 8 && (L_{2})\times(-2)}
\begin{calculs}
(12x+9y)+(-12x+10y) &= 30+8 &&\justif{on additionne}\\
19y &= 38 &&\justif{les $x$ s'éliminent}\\
y &= 2 &&\justif{on divise par $19$}
\end{calculs}
\begin{calculs}
4x+3\times 2 &= 10 &&\justif{dans $(L_{1})$}\\
4x &= 4 &&\justif{on enlève $6$}\\
x &= 1 &&\justif{on divise par $4$}
\end{calculs}
\explique{$S=\res{\{(1\,;2)\}}$.}{contrôle dans $(L_{2})$ : $6-10=-4$}
\cas{d)}
\explique{En $y$, les coefficients $3$ et $-3$ sont opposés : la combinaison est
immédiate.}{}
\begin{calculs}
(2x+3y)+(4x-3y) &= 2+1 &&\justif{$(L_{1})+(L_{2})$}\\
6x &= 3 &&\justif{les $y$ s'éliminent}\\
x &= \frac{3}{6} &&\justif{on divise par $6$}\\
x &= \frac{1}{2} &&\justif{on simplifie par $3$}
\end{calculs}
\begin{calculs}
2\times\frac{1}{2}+3y &= 2 &&\justif{dans $(L_{1})$}\\
1+3y &= 2 &&\justif{$2\times\frac12=1$}\\
3y &= 1 &&\justif{on enlève $1$}\\
y &= \frac{1}{3} &&\justif{on divise par $3$}
\end{calculs}
\explique{$S=\res{\left\{\left(\frac{1}{2}\,;\frac{1}{3}\right)\right\}}$.}{contrôle
dans $(L_{2})$ : $4\times\frac12-3\times\frac13=2-1=1$}
\end{exo}

\partie{Problèmes et systèmes}

\begin{exo}[Mangas et bandes dessinées]
\cas{1)}
\explique{Lou : un manga et deux bandes dessinées pour $38$~€, donc
$\res{x+2y=38}$.}{une information, une équation}
\explique{Sam : quatre mangas et trois bandes dessinées pour $77$~€, donc
$\res{4x+3y=77}$.}{}
\cas{2)}
\sysligne{x+2y &= 38 && (L_{1})\\ 4x+3y &= 77 && (L_{2})}
\begin{calculs}
x &= 38-2y &&\justif{on isole $x$ dans $(L_{1})$ : coefficient $1$}
\end{calculs}
\begin{calculs}
4x+3y &= 77 &&\justif{$(L_{2})$}\\
4(38-2y)+3y &= 77 &&\justif{on remplace $x$ par $38-2y$}\\
152-8y+3y &= 77 &&\justif{on développe}\\
-5y &= -75 &&\justif{on réduit et on enlève $152$}\\
y &= 15 &&\justif{on divise par $-5$}
\end{calculs}
\begin{calculs}
x &= 38-2\times 15 &&\justif{$x=38-2y$}\\
x &= 8 &&
\end{calculs}
\explique{$S=\res{\{(8\,;15)\}}$.}{contrôle dans $(L_{2})$ : $32+45=77$}
\cas{3)}
\explique{Un manga coûte \res{$8$~€} et une bande dessinée \res{$15$~€}.}{on répond
par une phrase}
\end{exo}

\begin{exo}[Imprimante 3D]
\explique{Soit $x$ le nombre de porte-clés et $y$ le nombre de figurines.}{on nomme les
inconnues}
\explique{$40$ objets, donc $x+y=40$ ; $20$~g par porte-clés et $50$~g par figurine, pour
$1\,220$~g en tout, donc $20x+50y=1\,220$.}{une information, une équation}
\sysligne{x+y &= 40 && (L_{1})\\ 20x+50y &= 1\,220 && (L_{2})}
\begin{calculs}
y &= 40-x &&\justif{on isole $y$ dans $(L_{1})$}
\end{calculs}
\begin{calculs}
20x+50y &= 1\,220 &&\justif{$(L_{2})$}\\
20x+50(40-x) &= 1\,220 &&\justif{on remplace $y$ par $40-x$}\\
20x+2\,000-50x &= 1\,220 &&\justif{on développe}\\
-30x &= -780 &&\justif{on réduit et on enlève $2\,000$}\\
x &= 26 &&\justif{on divise par $-30$}
\end{calculs}
\begin{calculs}
y &= 40-26 &&\justif{$y=40-x$}\\
y &= 14 &&
\end{calculs}
\explique{Le club a imprimé \res{$26$ porte-clés et $14$ figurines}.}{contrôle :
$520+700=1\,220$}
\end{exo}

\begin{exo}[Sortie scolaire]
\explique{Soit $x$ le nombre d'élèves et $y$ le nombre d'adultes.}{on nomme les
inconnues}
\explique{Chaque ligne du tableau donne une équation :}{prix par personne $\times$
nombre de personnes}
\sysligne{12x+20y &= 680 && (L_{1})\\ 10x+30y &= 620 && (L_{2})}
\explique{Aucun coefficient ne vaut $1$ : on combine pour éliminer $y$.}{$20\times 3=60$
et $30\times(-2)=-60$}
\sysligne{36x+60y &= 2\,040 && (L_{1})\times 3\\ -20x-60y &= -1\,240 && (L_{2})\times(-2)}
\begin{calculs}
(36x+60y)+(-20x-60y) &= 2\,040-1\,240 &&\justif{on additionne}\\
16x &= 800 &&\justif{les $y$ s'éliminent}\\
x &= 50 &&\justif{on divise par $16$}
\end{calculs}
\begin{calculs}
12\times 50+20y &= 680 &&\justif{dans $(L_{1})$}\\
600+20y &= 680 &&\\
20y &= 80 &&\justif{on enlève $600$}\\
y &= 4 &&\justif{on divise par $20$}
\end{calculs}
\explique{\res{$50$ élèves et $4$ adultes} participent à la sortie.}{contrôle avec B :
$500+120=620$}
\end{exo}

\partie{Systèmes et intersection de droites}

\begin{exo}[Point d'intersection]
\explique{On calcule le déterminant des vecteurs directeurs : non nul, les droites sont
sécantes ; puis on résout le système de leurs équations.}{méthode 4}
\cas{a)}
\explique{$\vec{u}_{1}\coord{-1}{1}$ dirige $(d_{1})$ et $\vec{u}_{2}\coord{1}{2}$ dirige
$(d_{2})$.}{$\vec{u}\coord{-b}{a}$ : pour $(d_{2})$, $-b=-(-1)=1$}
\begin{calculs}
\det(\vec{u}_{1}\,;\vec{u}_{2}) &= -1\times 2-1\times 1 &&\justif{$xy'-yx'$}\\
\det(\vec{u}_{1}\,;\vec{u}_{2}) &= -3 &&\justif{non nul : les droites sont sécantes}
\end{calculs}
\sysligne{x+y &= 5 && (L_{1})\\ 2x-y &= 1 && (L_{2})}
\begin{calculs}
(x+y)+(2x-y) &= 5+1 &&\justif{$(L_{1})+(L_{2})$ : coefficients de $y$ opposés}\\
3x &= 6 &&\\
x &= 2 &&\justif{on divise par $3$}
\end{calculs}
\explique{$y=5-2=3$ : le point d'intersection a pour coordonnées
$\res{(2\,;3)}$.}{dans $(L_{1})$ ; contrôle dans $(d_{2})$ : $4-3-1=0$}
\cas{b)}
\explique{$\vec{u}_{3}\coord{1}{3}$ et $\vec{u}_{4}\coord{1}{-1}$.}{$\coord{1}{m}$ pour
$y=mx+p$}
\begin{calculs}
\det(\vec{u}_{3}\,;\vec{u}_{4}) &= 1\times(-1)-3\times 1 &&\justif{$xy'-yx'$}\\
\det(\vec{u}_{3}\,;\vec{u}_{4}) &= -4 &&\justif{non nul : sécantes}
\end{calculs}
\begin{calculs}
3x-2 &= -x+6 &&\justif{même $y$ : on remplace $y$ dans la deuxième équation}\\
4x &= 8 &&\justif{on ajoute $x$ et $2$}\\
x &= 2 &&\justif{on divise par $4$}
\end{calculs}
\explique{$y=3\times 2-2=4$ : le point d'intersection a pour coordonnées
$\res{(2\,;4)}$.}{contrôle dans $(d_{4})$ : $-2+6=4$}
\cas{c)}
\explique{$\vec{u}_{5}\coord{-2}{3}$ et $\vec{u}_{6}\coord{4}{1}$.}{$\vec{u}\coord{-b}{a}$}
\begin{calculs}
\det(\vec{u}_{5}\,;\vec{u}_{6}) &= -2\times 1-3\times 4 &&\justif{$xy'-yx'$}\\
\det(\vec{u}_{5}\,;\vec{u}_{6}) &= -14 &&\justif{non nul : sécantes}
\end{calculs}
\sysligne{3x+2y &= 4 && (L_{1})\\ x-4y &= -8 && (L_{2})}
\begin{calculs}
x &= 4y-8 &&\justif{on isole $x$ dans $(L_{2})$ : on ajoute $4y$}
\end{calculs}
\begin{calculs}
3x+2y &= 4 &&\justif{$(L_{1})$}\\
3(4y-8)+2y &= 4 &&\justif{on remplace $x$ par $4y-8$}\\
14y-24 &= 4 &&\justif{on développe et on réduit}\\
14y &= 28 &&\justif{on ajoute $24$}\\
y &= 2 &&\justif{on divise par $14$}
\end{calculs}
\explique{$x=4\times 2-8=0$ : le point d'intersection a pour coordonnées
$\res{(0\,;2)}$.}{contrôle dans $(d_{5})$ : $0+4-4=0$}
\end{exo}

\begin{exo}[Parallèles ou sécantes ?]
\cas{a)}
\explique{$\vec{u}_{1}\coord{6}{2}$ dirige $(d_{1})$ et $\vec{u}_{2}\coord{-3}{-1}$ dirige
$(d_{2})$.}{$\vec{u}\coord{-b}{a}$}
\begin{calculs}
\det(\vec{u}_{1}\,;\vec{u}_{2}) &= 6\times(-1)-2\times(-3) &&\justif{$xy'-yx'$}\\
\det(\vec{u}_{1}\,;\vec{u}_{2}) &= 0 &&
\end{calculs}
\explique{Le déterminant est nul : les droites $(d_{1})$ et $(d_{2})$ sont
\res{parallèles}.}{vecteurs directeurs colinéaires}
\rem{Non exigé : en multipliant l'équation de $(d_{2})$ par $-2$, on obtient
$2x-6y-8=0$, qui n'est pas l'équation de $(d_{1})$ ; les droites sont strictement
parallèles (partie VI).}
\cas{b)}
\explique{$\vec{u}_{3}\coord{-1}{4}$ dirige $(d_{3})$ et $\vec{u}_{4}\coord{1}{2}$ dirige
$(d_{4})$.}{$\coord{1}{m}$ pour $y=2x$}
\begin{calculs}
\det(\vec{u}_{3}\,;\vec{u}_{4}) &= -1\times 2-4\times 1 &&\justif{$xy'-yx'$}\\
\det(\vec{u}_{3}\,;\vec{u}_{4}) &= -6 &&
\end{calculs}
\explique{Le déterminant n'est pas nul : les droites $(d_{3})$ et $(d_{4})$ sont
\res{sécantes}.}{}
\begin{calculs}
4x+y-6 &= 0 &&\justif{$(d_{3})$}\\
4x+2x-6 &= 0 &&\justif{on remplace $y$ par $2x$}\\
6x &= 6 &&\justif{on réduit et on ajoute $6$}\\
x &= 1 &&\justif{on divise par $6$}
\end{calculs}
\explique{$y=2\times 1=2$ : le point d'intersection a pour coordonnées
$\res{(1\,;2)}$.}{contrôle dans $(d_{3})$ : $4+2-6=0$}
\end{exo}

\begin{exo}[Deux droites à déterminer]
\cas{1)}
\explique{$A$ : $-1-2\times 1+3=0$ ; $B$ : $3-2\times 3+3=0$.}{on remplace $x$ et $y$
par les coordonnées}
\explique{Les points $A$ et $B$ vérifient l'équation $x-2y+3=0$ : c'est bien une
équation de la droite $(AB)$.}{deux points suffisent à fixer une droite}
\cas{2)}
\explique{$\vec{u}\coord{-1}{2}$ : $-b=-1$ et $a=2$, donc $b=1$. $(d')$ a une
équation de la forme $2x+y+c=0$.}{on lit $a$ et $b$ sur le vecteur directeur
(chapitre 9)}
\begin{calculs}
2\times 3+(-2)+c &= 0 &&\justif{$C(3\,;-2)$ est sur $(d')$}\\
4+c &= 0 &&\\
c &= -4 &&
\end{calculs}
\explique{$(d')$ : $\res{2x+y-4=0}$.}{}
\cas{3)}
\explique{$\vec{v}\coord{2}{1}$ dirige $(d)$ et $\vec{u}\coord{-1}{2}$ dirige
$(d')$.}{pour $(d)$ : $-b=-(-2)=2$}
\begin{calculs}
\det(\vec{v}\,;\vec{u}) &= 2\times 2-1\times(-1) &&\justif{$xy'-yx'$}\\
\det(\vec{v}\,;\vec{u}) &= 5 &&
\end{calculs}
\explique{Le déterminant n'est pas nul : $(d)$ et $(d')$ sont
\res{sécantes}.}{}
\cas{4)}
\sysligne{x-2y &= -3 && (L_{1})\\ 2x+y &= 4 && (L_{2})}
\begin{calculs}
x &= 2y-3 &&\justif{on isole $x$ dans $(L_{1})$}
\end{calculs}
\begin{calculs}
2x+y &= 4 &&\justif{$(L_{2})$}\\
2(2y-3)+y &= 4 &&\justif{on remplace $x$ par $2y-3$}\\
5y-6 &= 4 &&\justif{on développe et on réduit}\\
5y &= 10 &&\justif{on ajoute $6$}\\
y &= 2 &&\justif{on divise par $5$}
\end{calculs}
\explique{$x=2\times 2-3=1$, donc $K\res{(1\,;2)}$.}{contrôle dans $(L_{2})$ :
$2+2=4$}
\cas{5)}
\cotefig[0.4]{%
\explique{On trace $(d)$ par $A$ et $B$, puis $(d')$ par $C$ et $K$ (ou avec
$\vec{u}$ : un carreau vers la gauche, deux vers le haut).}{}
\explique{Les deux droites se coupent bien au point de coordonnées $(1\,;2)$.}{la
lecture contrôle le calcul}}{%
\begin{tikzpicture}[x=0.45cm,y=0.45cm]
  \repere{-2}{-3}{5}{4}
  \gradx{-1,1,2,3,4} \grady{-2,-1,1,2,3}
  \draw[droite] (-2,0.5) -- (5,4);
  \draw[droiteB] (0.5,3) -- (3.5,-3);
  \draw[aide] (1,0) -- (1,2) -- (0,2);
  \pt{-1,1}{above left}{A}
  \pt{3,3}{above left}{B}
  \pt{3,-2}{right}{C}
  \pt{1,2}{above left}{K}
  \node[coulCourbe,font=\footnotesize,below right] at (4,3.5) {$(d)$};
  \node[coulCourbeB,font=\footnotesize,left] at (2.2,-1) {$(d')$};
\end{tikzpicture}}
\end{exo}

\begin{exo}[Deux médianes d'un triangle]
\cas{1)}
\begin{calculs}
x_{I} &= \frac{2+(-4)}{2} = -1 &&\justif{milieu : moyenne des abscisses}\\
y_{I} &= \frac{6+0}{2} = 3 &&
\end{calculs}
\explique{$I\res{(-1\,;3)}$ ; de même, $J\res{(5\,;3)}$.}{$x_{J}=\frac{2+8}{2}=5$ et
$y_{J}=\frac{6+0}{2}=3$}
\cas{2)}
\explique{$C$ : $8+3\times 0-8=0$ ; $I$ : $-1+3\times 3-8=0$.}{on remplace par les
coordonnées}
\explique{$C$ et $I$ vérifient l'équation $x+3y-8=0$ : c'est une équation de
$(CI)$.}{}
\cas{3)}
\explique{$\vect{BJ}\coord{9}{3}$ ; on prend le vecteur directeur plus simple
$\vec{u}\coord{3}{1}$.}{$\vec{u}=\frac13\vect{BJ}$, colinéaire}
\explique{$-b=3$ et $a=1$ : $(BJ)$ a une équation de la forme $x-3y+c=0$.}{}
\begin{calculs}
-4-3\times 0+c &= 0 &&\justif{$B(-4\,;0)$ est sur $(BJ)$}\\
c &= 4 &&
\end{calculs}
\explique{$(BJ)$ : $\res{x-3y+4=0}$.}{contrôle avec $J$ : $5-9+4=0$}
\cas{4)}
\sysligne{x+3y &= 8 && (L_{1})\\ x-3y &= -4 && (L_{2})}
\begin{calculs}
(x+3y)+(x-3y) &= 8-4 &&\justif{$(L_{1})+(L_{2})$ : coefficients de $y$ opposés}\\
2x &= 4 &&\\
x &= 2 &&
\end{calculs}
\begin{calculs}
2+3y &= 8 &&\justif{dans $(L_{1})$}\\
3y &= 6 &&\\
y &= 2 &&
\end{calculs}
\explique{$G\res{(2\,;2)}$.}{contrôle dans $(L_{2})$ : $2-6=-4$}
\cas{5)}
\explique{$M\left(\frac{-4+8}{2}\,;\frac{0+0}{2}\right)$, soit $M(2\,;0)$. $A$ et $M$ ont
la même abscisse $2$, donc $(AM)$ est la droite verticale d'équation $x=2$.}{droites
parallèles aux axes (chapitre 9)}
\explique{$G$ a pour abscisse $2$, donc $G$ est sur $(AM)$.}{}
\explique{$G$ est sur les trois médianes du triangle $ABC$ : c'est son \res{centre de
gravité}.}{les trois médianes sont concourantes}
\end{exo}

\partie{Équation réduite d'une droite par un système}

\begin{exo}[Par deux points]
\explique{$(AB)$ a pour équation $y=mx+p$ ; chaque point donne une équation d'inconnues
$m$ et $p$.}{méthode 5}
\cas{a)}
\sysligne{m+p &= 4 && (L_{1})\\ 3m+p &= 10 && (L_{2})}
\begin{calculs}
(-m-p)+(3m+p) &= -4+10 &&\justif{$(L_{1})\times(-1)$, puis on additionne $(L_{2})$}\\
2m &= 6 &&\justif{les $p$ s'éliminent}\\
m &= 3 &&
\end{calculs}
\explique{$3+p=4$, donc $p=1$ : $(AB)$ a pour équation $\res{y=3x+1}$.}{dans
$(L_{1})$ ; contrôle avec $B$ : $9+1=10$}
\cas{b)}
\sysligne{-2m+p &= 7 && (L_{1})\\ 2m+p &= -1 && (L_{2})}
\begin{calculs}
(-2m+p)+(2m+p) &= 7-1 &&\justif{coefficients de $m$ déjà opposés}\\
2p &= 6 &&\justif{les $m$ s'éliminent}\\
p &= 3 &&
\end{calculs}
\begin{calculs}
2m+3 &= -1 &&\justif{dans $(L_{2})$}\\
2m &= -4 &&\\
m &= -2 &&
\end{calculs}
\explique{$(AB)$ : $\res{y=-2x+3}$.}{contrôle avec $A$ : $4+3=7$}
\cas{c)}
\sysligne{-m+p &= 0 && (L_{1})\\ 2m+p &= 1 && (L_{2})}
\begin{calculs}
(m-p)+(2m+p) &= 0+1 &&\justif{$(L_{1})\times(-1)$, puis on additionne $(L_{2})$}\\
3m &= 1 &&\\
m &= \frac{1}{3} &&
\end{calculs}
\explique{$(L_{1})$ donne $p=m=\frac13$ : $(AB)$ a pour équation
$\res{y=\frac{1}{3}x+\frac{1}{3}}$.}{contrôle avec $B$ : $\frac23+\frac13=1$}
\cas{d)}
\sysligne{2m+p &= -1 && (L_{1})\\ 5m+p &= 1 && (L_{2})}
\begin{calculs}
(-2m-p)+(5m+p) &= 1+1 &&\justif{$(L_{1})\times(-1)$, puis on additionne $(L_{2})$}\\
3m &= 2 &&\\
m &= \frac{2}{3} &&
\end{calculs}
\begin{calculs}
2\times\frac{2}{3}+p &= -1 &&\justif{dans $(L_{1})$}\\
p &= -\frac{3}{3}-\frac{4}{3} &&\justif{on enlève $\frac43$ ; $-1=-\frac33$}\\
p &= -\frac{7}{3} &&
\end{calculs}
\explique{$(AB)$ : $\res{y=\frac{2}{3}x-\frac{7}{3}}$.}{contrôle avec $B$ :
$\frac{10}{3}-\frac73=\frac33=1$}
\end{exo}

\begin{exo}[Deux méthodes, un résultat]
\cas{1)}
\sysligne{-3m+p &= 5 && (L_{1})\\ m+p &= -3 && (L_{2})}
\begin{calculs}
(3m-p)+(m+p) &= -5-3 &&\justif{$(L_{1})\times(-1)$, puis on additionne $(L_{2})$}\\
4m &= -8 &&\\
m &= -2 &&
\end{calculs}
\explique{$-2+p=-3$, donc $p=-1$ : $(AB)$ a pour équation $\res{y=-2x-1}$.}{dans
$(L_{2})$ ; contrôle avec $A$ : $6-1=5$}
\cas{2)}
\begin{calculs}
m &= \frac{y_{B}-y_{A}}{x_{B}-x_{A}} &&\justif{formule de la pente (chapitre 5)}\\
m &= \frac{-3-5}{1-(-3)} &&\\
m &= \frac{-8}{4} &&\\
m &= -2 &&
\end{calculs}
\explique{On retrouve bien $\res{m=-2}$.}{}
\cas{3)}
\explique{$-2\times 4-1=-9$ : les coordonnées de $C$ vérifient l'équation de $(AB)$,
donc $C$ est sur $(AB)$ ; les points $A$, $B$ et $C$ \res{sont alignés}.}{on remplace
$x$ par $x_{C}=4$}
\end{exo}

\begin{exo}[Salle de sport]
\cas{1)}
\explique{$10$ séances coûtent $95$~€, donc $95=m\times 10+p$ ; $25$ séances coûtent
$170$~€, donc $170=m\times 25+p$.}{le point $(10\,;95)$ est sur la droite}
\explique{Le système est $\res{\sys{10m+p &= 95 && (L_{1})\\ 25m+p &= 170 &&
(L_{2})}}$}{}
\cas{2)}
\begin{calculs}
(-10m-p)+(25m+p) &= -95+170 &&\justif{$(L_{1})\times(-1)$, puis on additionne $(L_{2})$}\\
15m &= 75 &&\justif{les $p$ s'éliminent}\\
m &= 5 &&\justif{on divise par $15$}
\end{calculs}
\begin{calculs}
10\times 5+p &= 95 &&\justif{dans $(L_{1})$}\\
p &= 45 &&\justif{on enlève $50$}
\end{calculs}
\explique{$\res{m=5}$ et $\res{p=45}$.}{contrôle dans $(L_{2})$ : $125+45=170$}
\cas{3)}
\explique{$m$ est le \res{prix d'une séance}, $5$~€ ; $p$ est le \res{prix de
l'abonnement}, $45$~€.}{$y$ augmente de $m$ à chaque séance ; pour $0$ séance,
$y=p$}
\cas{4)}
\begin{calculs}
y &= 5\times 40+45 &&\justif{$y=5x+45$ avec $x=40$}\\
y &= 245 &&
\end{calculs}
\explique{Pour $40$ séances, on paie \res{$245$~€}.}{}
\end{exo}

\partie{Aucune ou une infinité de solutions}

\begin{exo}[Combien de solutions ?]
\explique{Si les deux inconnues disparaissent : égalité fausse, aucune solution ;
égalité toujours vraie, une infinité.}{propriété 3, méthode 6}
\cas{a)}
\begin{calculs}
x &= 3+2y &&\justif{on isole $x$ dans $(L_{1})$}
\end{calculs}
\begin{calculs}
-3x+6y &= 2 &&\justif{$(L_{2})$}\\
-3(3+2y)+6y &= 2 &&\justif{on remplace $x$ par $3+2y$}\\
-9-6y+6y &= 2 &&\justif{on développe}\\
-9 &= 2 &&\justif{les $y$ ont disparu : égalité fausse}
\end{calculs}
\explique{$S=\res{\varnothing}$ : les deux droites sont strictement
parallèles.}{aucun point commun}
\cas{b)}
\begin{calculs}
y &= 4-2x &&\justif{on isole $y$ dans $(L_{1})$}
\end{calculs}
\begin{calculs}
6x+3y &= 12 &&\justif{$(L_{2})$}\\
6x+3(4-2x) &= 12 &&\justif{on remplace $y$ par $4-2x$}\\
6x+12-6x &= 12 &&\justif{on développe}\\
12 &= 12 &&\justif{égalité toujours vraie}
\end{calculs}
\explique{Le système a \res{une infinité de solutions} : les couples $(x\,;y)$ tels que
$2x+y=4$, par exemple $(0\,;4)$ et $(2\,;0)$. Les deux droites sont
confondues.}{$(L_{2})$ est $(L_{1})$ multipliée par $3$}
\cas{c)}
\begin{calculs}
y &= 5-x &&\justif{on isole $y$ dans $(L_{1})$}
\end{calculs}
\begin{calculs}
2x+y &= 7 &&\justif{$(L_{2})$}\\
2x+(5-x) &= 7 &&\justif{on remplace $y$ par $5-x$}\\
x+5 &= 7 &&\justif{on réduit}\\
x &= 2 &&\justif{on enlève $5$}
\end{calculs}
\explique{$y=5-2=3$, donc $S=\res{\{(2\,;3)\}}$ : les deux droites sont sécantes au
point de coordonnées $(2\,;3)$.}{une seule solution}
\cas{d)}
\begin{calculs}
4x-2y &= 6 &&\justif{$(L_{1})$}\\
-2y &= 6-4x &&\justif{on enlève $4x$}\\
y &= 2x-3 &&\justif{on divise par $-2$}
\end{calculs}
\begin{calculs}
-6x+3y &= -9 &&\justif{$(L_{2})$}\\
-6x+3(2x-3) &= -9 &&\justif{on remplace $y$ par $2x-3$}\\
-6x+6x-9 &= -9 &&\justif{on développe}\\
-9 &= -9 &&\justif{égalité toujours vraie}
\end{calculs}
\explique{Le système a \res{une infinité de solutions} : les couples $(x\,;y)$ tels que
$y=2x-3$, par exemple $(0\,;-3)$ et $(3\,;3)$. Les deux droites sont
confondues.}{les deux équations s'écrivent $y=2x-3$}
\rem{Erreur fréquente : conclure « $S=\{(0\,;-3)\}$ ». Il n'y a pas une solution, mais
une infinité.}
\end{exo}

\begin{exo}[Un coefficient inconnu]
\cas{1)}
\explique{$(d)$ : $\res{\vec{u}\coord{-1}{2}}$ ; $(d')$ :
$\res{\vec{u}'\coord{-3}{a}}$.}{$\vec{u}\coord{-b}{a}$ : pour $(d')$, le coefficient
de $x$ est $a$}
\cas{2)}
\begin{calculs}
\det(\vec{u}\,;\vec{u}') &= -1\times a-2\times(-3) &&\justif{$xy'-yx'$}\\
\det(\vec{u}\,;\vec{u}') &= \res{-a+6} &&
\end{calculs}
\cas{3)}
\begin{calculs}
-a+6 &= 0 &&\justif{parallèles : déterminant nul}\\
a &= 6 &&
\end{calculs}
\explique{Les droites sont parallèles pour \res{$a=6$}.}{}
\cas{4)}
\sysligne{2x+y &= 3 && (L_{1})\\ 6x+3y &= 1 && (L_{2})}
\begin{calculs}
y &= 3-2x &&\justif{on isole $y$ dans $(L_{1})$}
\end{calculs}
\begin{calculs}
6x+3(3-2x) &= 1 &&\justif{dans $(L_{2})$, on remplace $y$}\\
6x+9-6x &= 1 &&\justif{on développe}\\
9 &= 1 &&\justif{égalité fausse}
\end{calculs}
\explique{$S=\res{\varnothing}$ : pour $a=6$, les droites $(d)$ et $(d')$ sont
strictement parallèles.}{aucun point commun}
\cas{5)}
\explique{Pour $a=4$, le déterminant vaut $-4+6=2$ : les droites sont sécantes.}{}
\sysligne{2x+y &= 3 && (L_{1})\\ 4x+3y &= 1 && (L_{2})}
\begin{calculs}
4x+3(3-2x) &= 1 &&\justif{dans $(L_{2})$, on remplace $y$ par $3-2x$}\\
4x+9-6x &= 1 &&\justif{on développe}\\
-2x &= -8 &&\justif{on réduit et on enlève $9$}\\
x &= 4 &&\justif{on divise par $-2$}
\end{calculs}
\explique{$y=3-2\times 4=-5$ : le point d'intersection a pour coordonnées
$\res{(4\,;-5)}$.}{contrôle dans $(L_{2})$ : $16-15=1$}
\end{exo}

\begin{exo}[Le glacier et son perroquet]
\cas{1)}
\explique{Deux fois plus cher : $4$~€ la vanille et $6$~€ le chocolat.}{$2\times 2$ et
$2\times 3$}
\sysligne{2x+3y &= 150 && (L_{1})\\ 4x+6y &= 310 && (L_{2})}
\begin{calculs}
(-4x-6y)+(4x+6y) &= -300+310 &&\justif{$(L_{1})\times(-2)$, puis on additionne $(L_{2})$}\\
0 &= 10 &&\justif{les deux inconnues disparaissent : égalité fausse}
\end{calculs}
\explique{Le système n'a aucune solution : le perroquet \res{a raison}.}{en doublant
les prix, la recette double : $300$~€, et non $310$~€}
\cas{2)}
\begin{calculs}
(-4x-6y)+(4x+6y) &= -300+300 &&\justif{même calcul, avec $300$}\\
0 &= 0 &&\justif{égalité toujours vraie}
\end{calculs}
\explique{Le système a une infinité de solutions : le perroquet \res{ne peut pas
trouver} $x$ et $y$.}{$(L_{2})$ est $(L_{1})$ multipliée par $2$ : elle n'apporte
rien de nouveau}
\rem{Par exemple, $30$ glaces de chaque parfum, ou $75$ à la vanille et aucune au
chocolat, donnent toutes deux $150$~€.}
\cas{3)}
\sysligne{2x+3y &= 150 && (L_{1})\\ 3x+2y &= 160 && (L_{2})}
\sysligne{6x+9y &= 450 && (L_{1})\times 3\\ -6x-4y &= -320 && (L_{2})\times(-2)}
\begin{calculs}
(6x+9y)+(-6x-4y) &= 450-320 &&\justif{on additionne}\\
5y &= 130 &&\justif{les $x$ s'éliminent}\\
y &= 26 &&\justif{on divise par $5$}
\end{calculs}
\begin{calculs}
2x+3\times 26 &= 150 &&\justif{dans $(L_{1})$}\\
2x &= 72 &&\justif{on enlève $78$}\\
x &= 36 &&\justif{on divise par $2$}
\end{calculs}
\explique{Le glacier a vendu \res{$36$ glaces à la vanille et $26$ au
chocolat}.}{contrôle dans $(L_{2})$ : $108+52=160$}
\end{exo}

\end{document}
